% Standalone certified dossier for Lipschitz transport and finite-group quotient monotonicity.
\documentclass[../main.tex]{subfiles}
\begin{document}

% === SOLUTION HEADER (proof provenance) =====================================
% ledger-nodes: lem:a5-lipschitz; thm:a5-monotone
% refines: lem:a5-lipschitz; thm:a5-monotone
% bounded_by: none
% evidence_run: none
% checked_by: agent
% author: /root/ab_legacy_author
% reviewer: /root/kls_core_author
% review: research/reviews/2026-08-25-ab-inline-r2-audit.md
% date: 2026-08-25
% ============================================================================

\section*{Solution: Lipschitz transport and quotient monotonicity}
\label{sec:sol-a5-lipschitz-quotient}

\paragraph{Analytic setting.}
For the Poincar\'e and log-Sobolev statements, the Dirichlet forms are understood on their closed
Sobolev domains.  Thus the standard weak chain rule is available: if \(T:X\to Y\) is globally
\(L\)-Lipschitz and \(f\) is admissible on \(Y\), then
\begin{equation}\label{eq:sol-a5-chain}
 \abs{\nabla(f\circ T)}\le L\,\abs{\nabla f}\circ T
 \quad\text{almost everywhere}.
\end{equation}
This includes Euclidean and Riemannian spaces by Rademacher's theorem and closure, and the
finite-group orbit spaces below with the lifted quotient form.  All measures considered for
\(T_2\) lie in \(\mathcal P_2\) of their respective complete separable metric spaces.

\subsection*{1. Lipschitz images}

\begin{lemma}[Ledger node \texttt{lem:a5-lipschitz}]
\label{lem:sol-a5-lipschitz}
Let \(\nu\) be a probability measure on \(X\), let \(T:X\to Y\) be measurable and globally
\(L\)-Lipschitz, and put \(\mu=T_\#\nu\).  Under the analytic setting above,
\[
 \CP(\mu)\le L^2\CP(\nu),\qquad
 \CLS(\mu)\le L^2\CLS(\nu),\qquad
 \CTCI(\mu)\le L^2\CTCI(\nu).
\]
No injectivity or measurable inverse for \(T\) is required.
\end{lemma}

\begin{proof}
For every admissible \(f:Y\to\R\), the pushforward identity and
\eqref{eq:sol-a5-chain} give
\[
 \Var_\mu(f)=\Var_\nu(f\circ T)
 \le \CP(\nu)\int_X\abs{\nabla(f\circ T)}^2\dd\nu
 \le L^2\CP(\nu)\int_Y\abs{\nabla f}^2\dd\mu.
\]
The same calculation starts from
\(\Ent_\mu(f^2)=\Ent_\nu((f\circ T)^2)\) and proves the log-Sobolev bound, including the
factor two in the repository normalization.

For transport, let \(\eta\ll\mu\), put \(h=\dd\eta/\dd\mu\), and define a measure on \(X\) by
\begin{equation}\label{eq:sol-a5-density-lift}
 \dd\widetilde\eta=(h\circ T)\dd\nu.
\end{equation}
This is a probability measure, since
\[
 \int_Xh\circ T\dd\nu=\int_Yh\dd\mu=1.
\]
Moreover, for every bounded measurable \(\phi\),
\[
 \int_Y\phi\dd(T_\#\widetilde\eta)
 =\int_X(\phi h)\circ T\dd\nu
 =\int_Y\phi h\dd\mu=\int_Y\phi\dd\eta,
\]
so \(T_\#\widetilde\eta=\eta\), and
\begin{equation}\label{eq:sol-a5-entropy-lift}
 \KL(\widetilde\eta\|\nu)
 =\int_X(h\log h)\circ T\dd\nu
 =\int_Yh\log h\dd\mu
 =\KL(\eta\|\mu).
\end{equation}
If \(\Gamma\) is any coupling of \(\widetilde\eta\) and \(\nu\), then
\((T,T)_\#\Gamma\) couples \(\eta\) and \(\mu\), and its quadratic cost is at most
\(L^2\) times that of \(\Gamma\).  Taking infima over couplings and using \(T_2\) for \(\nu\)
therefore gives
\[
 W_{2,Y}^2(\eta,\mu)
 \le L^2W_{2,X}^2(\widetilde\eta,\nu)
 \le 2L^2\CTCI(\nu)\KL(\widetilde\eta\|\nu)
 =2L^2\CTCI(\nu)\KL(\eta\|\mu).
\]
The construction \eqref{eq:sol-a5-density-lift} is constant on fibers and is valid even when
\(T\) is noninjective; no choice of a preimage was made.
\end{proof}

\subsection*{2. Finite isometric quotients}

\begin{theorem}[Ledger node \texttt{thm:a5-monotone}]
\label{thm:sol-a5-monotone}
Let a finite group \(G\) act isometrically on \((\Theta,d_\Theta)\), let \(\pi\) be a
\(G\)-invariant probability measure, and let
\[
 q:\Theta\to\bar\Theta=\Theta/G,
 \qquad \bar\pi=q_\#\pi,
 \qquad
 d_{\bar\Theta}([x],[y])=\min_{g\in G}d_\Theta(x,gy).
\]
Equip \(\bar\Theta\) with the lifted quotient Dirichlet form from
Equation~\ref{eq:a5-quotient-form}.  Then
\[
 \CP(\bar\pi)\le\CP(\pi),\qquad
 \CLS(\bar\pi)\le\CLS(\pi),\qquad
 \CTCI(\bar\pi)\le\CTCI(\pi).
\]
\end{theorem}

\begin{proof}
First, \(\bar\pi\) really is a probability measure:
\(\bar\pi(\bar\Theta)=\pi(q^{-1}(\bar\Theta))=\pi(\Theta)=1\).  Pullback by \(q\) preserves
integrals, hence variance and entropy,
\[
 \Var_{\bar\pi}(\varphi)=\Var_\pi(\varphi\circ q),
 \qquad
 \Ent_{\bar\pi}(\varphi^2)=\Ent_\pi((\varphi\circ q)^2).
\]
By the definition of the quotient form,
\[
 \mathcal E_{\bar\pi}(\varphi)
 =\mathcal E_\pi(\varphi\circ q).
\]
Thus each quotient test function lifts to an admissible \(G\)-invariant test function upstairs
with exactly the same Rayleigh or entropy-energy quotient.  Applying the inequalities for
\(\pi\), or equivalently restricting their variational suprema to this invariant subclass,
proves the Poincar\'e and log-Sobolev comparisons.

For completeness, we spell out the transport descent rather than merely cite
Lemma~\ref{lem:sol-a5-lipschitz}.  The quotient map is \(1\)-Lipschitz because the identity
element is one candidate in the minimum:
\[
 d_{\bar\Theta}(q(x),q(y))\le d_\Theta(x,y).
\]
Given \(\bar\eta\ll\bar\pi\), set \(h=\dd\bar\eta/\dd\bar\pi\) and
\(\dd\eta=(h\circ q)\dd\pi\).  The same pushforward calculation as above proves, without any
choice of orbit representative,
\[
 \eta(\Theta)=1,
 \qquad q_\#\eta=\bar\eta,
 \qquad \KL(\eta\|\pi)=\KL(\bar\eta\|\bar\pi).
\]
If \(\Gamma\) couples \(\eta\) and \(\pi\), then \((q,q)_\#\Gamma\) couples
\(\bar\eta\) and \(\bar\pi\), and the quotient cost is no larger than the upstairs cost.
Consequently
\[
 W_{2,\bar\Theta}^2(\bar\eta,\bar\pi)
 \le W_{2,\Theta}^2(\eta,\pi)
 \le2\CTCI(\pi)\KL(\eta\|\pi)
 =2\CTCI(\pi)\KL(\bar\eta\|\bar\pi),
\]
which proves \(\CTCI(\bar\pi)\le\CTCI(\pi)\).
\end{proof}

\paragraph{Scope.}
Monotonicity makes no assertion of strict improvement.  Invariant physical slow modes can survive
the quotient unchanged; their quantitative separation from label-switching modes is the subject
of the later invariant/non-invariant gap results.

\end{document}
